\documentclass[10pt,a4paper]{amsart}
\usepackage[margin=19mm]{geometry}
\usepackage{amsmath,amssymb,mathtools}
\usepackage[scr=rsfs,cal=euler]{mathalfa}
\usepackage{xfrac}
\usepackage[unicode,hidelinks,bookmarks=false]{hyperref}
\newcommand{\ie}{i.e.\ }
\newcommand{\cf}{cf.\ }
\newcommand{\resp}{respectively\ }
\newcommand{\Subsection}{\subsection}
\providecommand{\nobreakdash}{\nobreak\hbox{-}}
\setlength{\emergencystretch}{2em}
\multlinegap0pt
\usepackage{amssymb}
\newtheorem{corollary}{Corollary}
\newtheorem{lemma}{Lemma}
\newtheorem{proposition}{Proposition}
\newtheorem{theorem}{Theorem}
\title{Noncommutative localizations of a contractive quantum plane}
\author{Anar Dosi}
\date{Source: arXiv:2509.23539v2 (CC BY 4.0)}
\begin{document}
\maketitle

\noindent\emph{Source and licence.} Noncommutative localizations of a contractive quantum plane. Version arXiv:2509.23539v2 (CC BY 4.0), distributed by arXiv under the Creative Commons Attribution 4.0 International licence, http://creativecommons.org/licenses/by/4.0/, as recorded in the arXiv licence metadata for this exact version and shown on the abstract page https://arxiv.org/abs/2509.23539v2. Full licence notice: \texttt{licenses/arxiv-2509.23539-CC-BY-4.0.txt}.\par\smallskip

\noindent\textit{Editorial scope.} Counted unit: the article's lemma lemOMN (source0.tex lines 5051--5071). The article's complete original proof of it is delivered with it (source0.tex lines 5073--5281, one \texttt{proof} environment of 209 lines). No mathematics is written, repaired or re-derived: the statement and the proof are the article's own lines, in the article's own notation, and no retained line is substituted or re-flowed.  Retained uncounted as the source-local context the proof needs: lines 3801--3816; lines 3834--3860; lines 4146--4151; lines 4478--4494; lines 4552--4557; lines 4703--4780.  Nothing was omitted from any retained span, so the package carries no omission record.  No retained line contains a citation, so the wrapper carries no bibliography and no bibliography entry.  No retained line contains a figure environment, an image-inclusion command or drawing code, so no image is delivered and the package declares no asset dependency.  Editorial formatting only, in addition: an AMS \texttt{amsart} preamble that restates the article's own declarations and macros used by the retained lines, namely the environments \texttt{corollary}, \texttt{lemma}, \texttt{proposition}, \texttt{theorem} on separate counters, together with the standard packages those lines need.  The retained environments print in this document as Corollary 1, Lemma 1, Proposition 1, Theorem 1 in order of appearance, whereas the article prints the same items with the numbering of its own full document.  The complete native source of the article as distributed in the arXiv e-print for this exact version is bundled unchanged as \texttt{sources/185893/source0.tex}.
\begin{corollary}
\label{corTech1}If $\zeta\in\mathcal{O}_{d,l}$ is a compatible quadruple then
\begin{align*}
z_{1}\left(  \zeta\right)   &  =\left(  \left(  z_{1}\zeta_{1}\left(
z_{1},z_{2}\right)  ,z_{1}\zeta_{2}\left(  z_{1},w_{2}\right)  \right)
,\left(  0,0\right)  \right)  ,\quad z_{2}\left(  \zeta\right)  =\left(
\left(  z_{2}\zeta_{1}\left(  z_{1},z_{2}\right)  ,0\right)  ,\left(
z_{2}\zeta_{3}\left(  w_{1},z_{2}\right)  ,0\right)  \right)  ,\\
w_{1}\left(  \zeta\right)   &  =\left(  \left(  0,0\right)  ,\left(
w_{1}\zeta_{3}\left(  w_{1},z_{2}\right)  ,w_{1}\zeta_{4}\left(  w_{1}%
,w_{2}\right)  \right)  \right)  ,\quad w_{2}\left(  \zeta\right)  =\left(
\left(  0,w_{2}\zeta_{2}\left(  z_{1},w_{2}\right)  \right)  ,\left(
0,w_{2}\zeta_{4}\left(  w_{1},w_{2}\right)  \right)  \right)  .
\end{align*}

\end{corollary}

\begin{lemma}
\label{lemTech3}If $\zeta\in\mathcal{O}_{d,l}$ is a compatible quadruple, then
the following identities
\begin{align*}
M_{x_{2},l}\left(  \zeta\right)   &  =\left(  \left(  \frac{\partial\zeta_{2}%
}{\partial w_{2}}\left(  z_{1},0\right)  ,\left(  \zeta_{2}\right)  _{w_{2}%
}\left(  z_{1},w_{2}\right)  \right)  ,\left(  \frac{\partial\zeta_{4}%
}{\partial w_{2}}\left(  w_{1},0\right)  ,\left(  \zeta_{4}\right)  _{w_{2}%
}\left(  w_{1},w_{2}\right)  \right)  \right)  ,\\
M_{y_{2},l}\left(  \zeta\right)   &  =q^{l+1}\left(  \left(  \left(  \zeta
_{1}\right)  _{z_{2}}\left(  z_{1},qz_{2}\right)  ,\frac{\partial\zeta_{1}%
}{\partial z_{2}}\left(  z_{1},0\right)  \right)  ,\left(  \left(  \zeta
_{3}\right)  _{z_{2}}\left(  w_{1},qz_{2}\right)  ,\frac{\partial\zeta_{3}%
}{\partial z_{2}}\left(  w_{1},0\right)  \right)  \right)  ,\\
N_{x_{1},d}\left(  \zeta\right)   &  =\left(  \left(  \left(  \zeta
_{1}\right)  _{z_{1}}\left(  z_{1},z_{2}\right)  ,\left(  \zeta_{2}\right)
_{z_{1}}\left(  z_{1},w_{2}\right)  \right)  ,\left(  \frac{\partial\zeta_{1}%
}{\partial z_{1}}\left(  0,z_{2}\right)  ,\frac{\partial\zeta_{2}}{\partial
z_{1}}\left(  0,w_{2}\right)  \right)  \right)  ,\\
N_{y_{1},d}\left(  \zeta\right)   &  =q^{d+1}\left(  \left(  \frac
{\partial\zeta_{3}}{\partial w_{1}}\left(  0,z_{2}\right)  ,\frac
{\partial\zeta_{4}}{\partial w_{1}}\left(  0,w_{2}\right)  \right)  ,\left(
\left(  \zeta_{3}\right)  _{w_{1}}\left(  qw_{1},z_{2}\right)  ,\left(
\zeta_{4}\right)  _{w_{1}}\left(  qw_{1},w_{2}\right)  \right)  \right)
\end{align*}
hold.
\end{lemma}

\begin{lemma}
\label{lemTech4}The range of the operator $T$ equals to $\ker\left(
\partial_{d,l}^{1}\right)  $, and $T$ implements a topological isomorphism
$\mathcal{Z}\left(  U\right)  \overset{\simeq}{\longrightarrow}\ker\left(
\partial_{d,l}^{1}\right)  $ of the Fr\'{e}chet spaces.
\end{lemma}

\begin{proposition}
\label{propTechE}The canonical identification $H_{d,l}^{1}=\mathcal{O}\left(
U\right)  $ holds up to a topological isomorphism. In this case, $\ker\left(
\partial_{d,l}^{1}\right)  =\operatorname{im}\left(  \partial_{d,l}%
^{1}\right)  \oplus\mathcal{O}\left(  U\right)  $ is a topological direct sum
of the Fr\'{e}chet spaces given by a continuous projection $p_{d,l}%
\in\mathcal{L}\left(  \ker\left(  \partial_{d,l}^{1}\right)  \right)  $ whose
image being identified with $\mathcal{O}\left(  U\right)  $ consists of the
following free quadruples
\[
\operatorname{im}\left(  p_{d,l}\right)  =\left\{  \left(  f\left(
z_{1}\right)  ,0,\lambda,g\left(  w_{2}\right)  \right)  :f\in\mathcal{O}%
\left(  z_{1}\right)  ,\lambda\in\mathbb{C},g\in\mathcal{O}\left(
w_{2}\right)  \right\}  \subseteq\mathcal{Z}\left(  U\right)  .
\]

\end{proposition}

\begin{theorem}
\label{thMain1}If $H_{d,l}^{p}$, $0\leq p\leq3$ are the cohomology groups of
the diagonal complex $\mathcal{O}_{q}\left(  d,l\right)  $, then $H_{d,l}%
^{0}=\left\{  0\right\}  $, $H_{d,l}^{1}=\mathcal{O}\left(  U\right)  $,
$H_{d,l}^{2}=\left\{  0\right\}  $ and $H_{d,l}^{3}=\left\{  0\right\}  $.
\end{theorem}

\begin{corollary}
\label{corDecom}The identities $\pi_{d,l}\tau_{d,l}^{0}=1$ and $\tau_{d,l}%
^{0}\pi_{d,l}+\partial_{d,l}^{1}\tau_{d,l}^{1}=1$ hold. In this case,
$p_{d,l}^{0}=\tau_{d,l}^{0}\pi_{d,l}\in\mathcal{L}\left(  \mathcal{O}%
_{d,l}\right)  $ and $p_{d,l}^{1}=\tau_{d,l}^{1}\partial_{d,l}^{1}%
\in\mathcal{L}\left(  \mathcal{O}_{d,l}^{\oplus2}\right)  $ are the continuous
projections such that the complex $\mathcal{O}_{q}\left(  d,l\right)  $ is
decomposed in the following way
\[%
\begin{array}
[c]{ccccccc}%
0\rightarrow\mathcal{O}_{d,l} & \overset{\partial_{d,l}^{0}}{\longrightarrow}
& \mathcal{O}_{d,l}^{\oplus2} & \overset{\partial_{d,l}^{1}}{\longrightarrow}
& \mathcal{O}_{d,l} & \overset{\pi_{d,l}}{\longrightarrow} & \mathcal{O}%
_{d+l}\rightarrow0\\
&  & \parallel &  & \parallel &  & \\%
\begin{array}
[c]{c}%
\\
\\
\\
\\
0\rightarrow\mathcal{O}_{d,l}%
\end{array}
&
\begin{array}
[c]{c}%
\\
\\
\\
\\
\overset{\partial_{d,l}^{0}}{\longleftrightarrow}%
\end{array}
&
\begin{array}
[c]{c}%
\operatorname{im}\left(  p_{d,l}^{1}\right) \\
\oplus\\
\operatorname{im}\left(  p_{d,l}\right) \\
\oplus\\
\operatorname{im}\left(  \partial_{d,l}^{0}\right)
\end{array}
&
\begin{array}
[c]{c}%
\\
_{\tau_{d,l}^{1}}\nwarrow\searrow^{\partial_{d,l}^{1}}\\
\\
\\
\end{array}
&
\begin{array}
[c]{c}%
\operatorname{im}\left(  p_{d,l}^{0}\right) \\
\oplus\\
\operatorname{im}\left(  \partial_{d,l}^{1}\right) \\
\\
\end{array}
&
\begin{array}
[c]{c}%
\\
_{\tau_{d,l}^{0}}\nwarrow\searrow^{\pi_{d,l}}\\
\\
\\
\end{array}
&
\begin{array}
[c]{c}%
\\
\\
\mathcal{O}_{d+l}\rightarrow0\\
\\
\end{array}
\end{array}
\]

\end{corollary}

\begin{lemma}
\label{lemOMN}Let $\left\{  s_{i,j}^{\left(  k\right)  }\right\}
\subseteq\operatorname{im}\left(  p_{i,j}\right)  $ and $\left\{  \omega
_{i,j}^{\left(  k\right)  }\right\}  \subseteq\mathcal{O}_{i,j}^{\oplus2}$ be
sequences such that the limit
\[
\chi_{i,j}=\lim_{k}\partial_{i,j}^{1}\omega_{i,j}^{\left(  k\right)
}+M_{i,j-1}s_{i,j-1}^{\left(  k\right)  }-N_{i-1,j}s_{i-1,j}^{\left(
k\right)  }\quad\text{in\quad}\mathcal{O}_{i,j}%
\]
converges to some $\chi_{i,j}\in\operatorname{im}\left(  p_{i,j}^{0}\right)  $
for all $i$, $j$, $i+j\leq n$. Then there are limits $\lim_{k}s_{i,j}^{\left(
k\right)  }=s_{i,j}$, $i+j<n$ and $\lim_{k}\partial_{i,j}^{1}\omega
_{i,j}^{\left(  k\right)  }=\partial_{i,j}^{1}\omega_{i,j}$, $i+j\leq n$ such
that
\[
\chi_{i,j}=\partial_{i,j}^{1}\omega_{i,j}+M_{i,j-1}s_{i,j-1}-N_{i-1,j}%
s_{i-1,j}%
\]
holds for all $i,j$, $i+j\leq n$.
\end{lemma}

\begin{proof}
First prove that $\lim_{k}s_{i,j-1}^{\left(  k\right)  }=s_{i,j-1}$ converges
iff so is $\lim_{k}s_{i-1,j}^{\left(  k\right)  }=s_{i-1,j}$ whenever $i+j\leq
n$. In this case, we would have
\[
\chi_{i,j}-M_{i,j-1}s_{i,j-1}+N_{i-1,j}s_{i-1,j}=\lim_{k}\partial_{i,j}%
^{1}\omega_{i,j}^{\left(  k\right)  }%
\]
out of the continuity of the operators $M_{i,j-1}$ and $N_{i-1,j}$. Taking
into account that $\operatorname{im}\left(  \partial_{i,j}^{1}\right)  $ is
closed (see Theorem \ref{thMain1}), we obtain that $\lim_{k}\partial_{i,j}%
^{1}\omega_{i,j}^{\left(  k\right)  }=\partial_{i,j}^{1}\omega_{i,j}$ for some
$\omega_{i,j}$, and the result $\chi_{i,j}=\partial_{i,j}^{1}\omega
_{i,j}+M_{i,j-1}s_{i,j-1}-N_{i-1,j}s_{i-1,j}$ follows.

Now by assumption, we have
\begin{align*}
s_{i,j-1}^{\left(  k\right)  }  &  =p_{i,j-1}\left(  s_{i,j-1}^{\left(
k\right)  }\right)  =\left(  f^{\left(  k\right)  }\left(  z_{1}\right)
,0,\lambda^{\left(  k\right)  },g^{\left(  k\right)  }\left(  w_{2}\right)
\right)  ,\\
s_{i-1,j}^{\left(  k\right)  }  &  =p_{i-1,j}\left(  s_{i-1,j}^{\left(
k\right)  }\right)  =\left(  u^{\left(  k\right)  }\left(  z_{1}\right)
,0,\mu^{\left(  k\right)  },v^{\left(  k\right)  }\left(  w_{2}\right)
\right)
\end{align*}
are free quadruples thanks to Proposition \ref{propTechE}, where $\left\{
f^{\left(  k\right)  }\right\}  \subseteq\mathcal{O}\left(  z\right)  $,
$\left\{  g^{\left(  k\right)  }\right\}  \subseteq\mathcal{O}\left(
w\right)  $ and $\left\{  \lambda^{\left(  k\right)  }\right\}  \subseteq
\mathbb{C}$. By Lemma \ref{lemTech4}, these free quadruples are in turn
identified with the following compatible ones
\begin{align*}
T\left(  s_{i,j-1}^{\left(  k\right)  }\right)   &  =\left[
\begin{array}
[c]{c}%
\zeta_{i,j-1}^{\left(  k\right)  }\\
\eta_{i,j-1}^{\left(  k\right)  }%
\end{array}
\right]  =\left[
\begin{array}
[c]{c}%
\left(  \left(  0,0\right)  ,\left(  w_{1}\lambda^{\left(  k\right)  }%
,\lambda^{\left(  k\right)  }w_{1}+w_{1}w_{2}g^{\left(  k\right)  }\left(
w_{2}\right)  \right)  \right) \\
\left(  \left(  \lambda^{\left(  k\right)  }z_{2}+z_{1}z_{2}f^{\left(
k\right)  }\left(  z_{1}\right)  ,0\right)  ,\left(  \lambda^{\left(
k\right)  }z_{2},0\right)  \right)
\end{array}
\right]  \in\ker\left(  \partial_{i,j-1}^{1}\right) \\
T\left(  s_{i-1,j}^{\left(  k\right)  }\right)   &  =\left[
\begin{array}
[c]{c}%
\zeta_{i-1,j}^{\left(  k\right)  }\\
\eta_{i-1,j}^{\left(  k\right)  }%
\end{array}
\right]  =\left[
\begin{array}
[c]{c}%
\left(  \left(  0,0\right)  ,\left(  \mu^{\left(  k\right)  }w_{1}%
,\mu^{\left(  k\right)  }w_{1}+w_{1}w_{2}v^{\left(  k\right)  }\left(
w_{2}\right)  \right)  \right) \\
\left(  \left(  \mu^{\left(  k\right)  }z_{2}+z_{1}z_{2}u^{\left(  k\right)
}\left(  z_{1}\right)  ,0\right)  ,\left(  \mu^{\left(  k\right)  }%
z_{2},0\right)  \right)
\end{array}
\right]  \in\ker\left(  \partial_{i-1,j}^{1}\right)
\end{align*}
in $\mathcal{O}_{i,j-1}$ and $\mathcal{O}_{i-1,j}$, respectively. Moreover,
based on Corollary \ref{corDecom}, we have
\[
\chi_{i,j}=p_{i,j}^{0}\left(  \chi_{i,j}\right)  =\tau_{i,j}^{0}\left(
f_{ij},g_{ij}\right)  =\left(  \left(  f_{ij}\left(  z_{1}\right)
,-\lambda_{ij}+f_{ij}\left(  z_{1}\right)  +g_{ij}\left(  w_{2}\right)
\right)  ,\left(  \lambda_{ij},g_{ij}\left(  w_{2}\right)  \right)  \right)
\]
in $\mathcal{O}_{i,j}$ for some $\left(  f_{i,j},g_{i,j}\right)
\in\mathcal{O}_{i+j}$ with $f_{ij}\left(  0\right)  =\lambda_{ij}%
=g_{ij}\left(  0\right)  $. Using Lemma \ref{lemTech3}, we deduce that
\begin{align*}
M_{i,j-1}s_{i,j-1}^{\left(  k\right)  }  &  =M_{x_{2},j-1}\zeta_{i,j-1}%
^{\left(  k\right)  }+M_{y_{2},j-1}\eta_{i,j-1}^{\left(  k\right)  }=\left(
\left(  0,0\right)  ,\left(  w_{1}g^{\left(  k\right)  }\left(  0\right)
,w_{1}g^{\left(  k\right)  }\left(  w_{2}\right)  \right)  \right) \\
&  +q^{j}\left(  \left(  \lambda^{\left(  k\right)  }+z_{1}f^{\left(
k\right)  }\left(  z_{1}\right)  ,\lambda^{\left(  k\right)  }+z_{1}f^{\left(
k\right)  }\left(  z_{1}\right)  \right)  ,\left(  \lambda^{\left(  k\right)
},\lambda^{\left(  k\right)  }\right)  \right)
\end{align*}%
\[
=\left(  \left(  q^{j}\left(  \lambda^{\left(  k\right)  }+z_{1}f^{\left(
k\right)  }\left(  z_{1}\right)  \right)  ,q^{j}\left(  \lambda^{\left(
k\right)  }+z_{1}f^{\left(  k\right)  }\left(  z_{1}\right)  \right)  \right)
,\left(  q^{j}\lambda^{\left(  k\right)  }+w_{1}g^{\left(  k\right)  }\left(
0\right)  ,q^{j}\lambda^{\left(  k\right)  }+w_{1}g^{\left(  k\right)
}\left(  w_{2}\right)  \right)  \right)  .
\]
In a similar way, we deduce that
\begin{align*}
N_{i-1,j}s_{i-1,j}^{\left(  k\right)  }  &  =N_{y_{1},i-1}\zeta_{i-1,j}%
^{\left(  k\right)  }+N_{x_{1},i-1}\eta_{i-1,j}^{\left(  k\right)  }%
=q^{i}\left(  \left(  \mu^{\left(  k\right)  },\mu^{\left(  k\right)  }%
+w_{2}v^{\left(  k\right)  }\left(  w_{2}\right)  \right)  ,\left(
\mu^{\left(  k\right)  },\mu^{\left(  k\right)  }+w_{2}v^{\left(  k\right)
}\left(  w_{2}\right)  \right)  \right) \\
&  +\left(  \left(  z_{2}u^{\left(  k\right)  }\left(  z_{1}\right)
,0\right)  ,\left(  z_{2}u^{\left(  k\right)  }\left(  0\right)  ,0\right)
\right)
\end{align*}%
\[
=\left(  \left(  q^{i}\mu^{\left(  k\right)  }+z_{2}u^{\left(  k\right)
}\left(  z_{1}\right)  ,q^{i}\left(  \mu^{\left(  k\right)  }+w_{2}v^{\left(
k\right)  }\left(  w_{2}\right)  \right)  \right)  ,\left(  q^{i}\mu^{\left(
k\right)  }+z_{2}u^{\left(  k\right)  }\left(  0\right)  ,q^{i}\left(
\mu^{\left(  k\right)  }+w_{2}v^{\left(  k\right)  }\left(  w_{2}\right)
\right)  \right)  \right)  .
\]
If $\omega_{i,j}^{\left(  k\right)  }=\left[
\begin{array}
[c]{c}%
a^{\left(  k\right)  }\\
b^{\left(  k\right)  }%
\end{array}
\right]  \in\mathcal{O}_{i,j}^{\oplus2}$ for some compactible quadruples
$a^{\left(  k\right)  }$ and $b^{\left(  k\right)  }$, then $\partial
_{i,j}^{1}\omega_{i,j}^{\left(  k\right)  }=\left(  z_{2}-q^{i}z_{1}\right)
a^{\left(  k\right)  }+\left(  q^{j}w_{2}-w_{1}\right)  b^{\left(  k\right)
}$. Using Corollary \ref{corTech1}, we obtain that
\[
\partial_{i,j}^{1}\omega_{i,j}^{\left(  k\right)  }=\left(  \left(  \left(
z_{2}-q^{i}z_{1}\right)  a_{1}^{\left(  k\right)  },-q^{i}z_{1}a_{2}^{\left(
k\right)  }+q^{j}w_{2}b_{2}^{\left(  k\right)  }\right)  ,\left(  z_{2}%
a_{3}^{\left(  k\right)  }-w_{1}b_{3}^{\left(  k\right)  },\left(  q^{j}%
w_{2}-w_{1}\right)  b_{4}^{\left(  k\right)  }\right)  \right)  .
\]
Put $\chi^{\left(  k\right)  }=\partial_{i,j}^{1}\omega_{i,j}^{\left(
k\right)  }+M_{i,j-1}s_{i,j-1}^{\left(  k\right)  }-N_{i-1,j}s_{i-1,j}%
^{\left(  k\right)  }$ to be convergent to $\chi_{i,j}$. By adding up all
those compatible quadruples obtained above, we have
\begin{align*}
\chi_{1}^{\left(  k\right)  }  &  =\left(  z_{2}-q^{i}z_{1}\right)
a_{1}^{\left(  k\right)  }+q^{j}\left(  \lambda^{\left(  k\right)  }%
+z_{1}f^{\left(  k\right)  }\left(  z_{1}\right)  \right)  -q^{i}\mu^{\left(
k\right)  }-z_{2}u^{\left(  k\right)  }\left(  z_{1}\right)  \rightarrow
f_{ij}\left(  z_{1}\right)  ,\\
\chi_{2}^{\left(  k\right)  }  &  =-q^{i}z_{1}a_{2}^{\left(  k\right)  }%
+q^{j}w_{2}b_{2}^{\left(  k\right)  }+q^{j}\left(  \lambda^{\left(  k\right)
}+z_{1}f^{\left(  k\right)  }\left(  z_{1}\right)  \right)  -q^{i}\left(
\mu^{\left(  k\right)  }+w_{2}v^{\left(  k\right)  }\left(  w_{2}\right)
\right) \\
&  \rightarrow-\lambda_{ij}+f_{ij}\left(  z_{1}\right)  +g_{ij}\left(
w_{2}\right)  ,\\
\chi_{3}^{\left(  k\right)  }  &  =z_{2}a_{3}^{\left(  k\right)  }-w_{1}%
b_{3}^{\left(  k\right)  }+q^{j}\lambda^{\left(  k\right)  }+w_{1}g^{\left(
k\right)  }\left(  0\right)  -q^{i}\mu^{\left(  k\right)  }-z_{2}u^{\left(
k\right)  }\left(  0\right)  \rightarrow\lambda_{ij},\\
\chi_{4}^{\left(  k\right)  }  &  =\left(  q^{j}w_{2}-w_{1}\right)
b_{4}^{\left(  k\right)  }+q^{j}\lambda^{\left(  k\right)  }+w_{1}g^{\left(
k\right)  }\left(  w_{2}\right)  -q^{i}\left(  \mu^{\left(  k\right)  }%
+w_{2}v^{\left(  k\right)  }\left(  w_{2}\right)  \right)  \rightarrow
g_{ij}\left(  w_{2}\right)  .
\end{align*}
For $w_{1}=0$, $z_{2}=0$ in $\chi_{3}^{\left(  k\right)  }\rightarrow
\lambda_{ij}$, we derive that $q^{j}\lambda^{\left(  k\right)  }-q^{i}%
\mu^{\left(  k\right)  }\rightarrow\lambda_{ij}$, Along the line $z_{2}%
=q^{i}z_{1}$ in $\chi_{1}^{\left(  k\right)  }\rightarrow f_{ij}\left(
z_{1}\right)  $, we obtain that $q^{j}f^{\left(  k\right)  }\left(
z_{1}\right)  -q^{i}u^{\left(  k\right)  }\left(  z_{1}\right)  \rightarrow
\left(  f_{ij}\right)  _{z}\left(  z_{1}\right)  $. Along the line
$w_{1}=q^{j}w_{2}$ in $\chi_{4}^{\left(  k\right)  }\rightarrow g_{ij}\left(
w_{2}\right)  $, we obtain that $q^{j}g^{\left(  k\right)  }\left(
w_{2}\right)  -q^{i}v^{\left(  k\right)  }\left(  w_{2}\right)  \rightarrow
\left(  g_{ij}\right)  _{w}\left(  w_{2}\right)  $. It follows that
\begin{align*}
q^{j}s_{i,j-1}^{\left(  k\right)  }-q^{i}s_{i-1,j}^{\left(  k\right)  }  &
=\left(  q^{j}f^{\left(  k\right)  }\left(  z_{1}\right)  -q^{i}u^{\left(
k\right)  }\left(  z_{1}\right)  ,0,q^{j}\lambda^{\left(  k\right)  }-q^{i}%
\mu^{\left(  k\right)  },q^{j}g^{\left(  k\right)  }\left(  w_{2}\right)
-q^{i}v^{\left(  k\right)  }\left(  w_{2}\right)  \right) \\
&  \rightarrow\left(  \left(  f_{ij}\right)  _{z}\left(  z_{1}\right)
,0,\lambda_{ij},\left(  g_{ij}\right)  _{w}\left(  w_{2}\right)  \right)  ,
\end{align*}
that is, the presence of $\lim_{k}s_{i,j-1}^{\left(  k\right)  }=s_{i,j-1}$
implies $\lim_{k}s_{i-1,j}^{\left(  k\right)  }=s_{i-1,j}$ and vise-versa. In
particular, if $j=0$ then the limit $\lim_{k}s_{i-1,0}^{\left(  k\right)
}=s_{i-1,0}$ ($s_{i,-1}^{\left(  k\right)  }=0$) does exist for every $i$. In
a similar way, $\lim_{k}s_{0,j-1}^{\left(  k\right)  }=s_{0,j-1}$
($s_{-1,j}^{\left(  k\right)  }=0$) exists for every $j$. Thus $s_{i,0}$,
$s_{0,j}$ do exist, and $\chi_{i,0}=\partial_{i,0}^{1}\omega_{i,0}%
-N_{i-1,0}s_{i-1,0}$, $\chi_{0,j}=\partial_{0,j}^{1}\omega_{0,j}%
+M_{0,j-1}s_{0,j-1}$ hold.

Finally, fix $m\leq n$. Then $s_{m-1,0}=\lim_{k}s_{m-1,0}^{\left(  k\right)
}$ implies that $\lim_{k}s_{m-2,1}^{\left(  k\right)  }=s_{m-2,1}$ converges
too. By iterating, we deduce that all $\lim_{k}s_{i,j}^{\left(  k\right)
}=s_{i,j}$, $i+j=m-1$ converge and
\begin{align*}
\chi_{m-i,i}  &  =\lim_{k}\partial_{m-i,i}^{1}\omega_{m-i,i}^{\left(
k\right)  }+M_{m-i,i-1}s_{m-i,i-1}^{\left(  k\right)  }-N_{m-i-1,i}%
s_{m-i-1,i}^{\left(  k\right)  }\\
&  =\partial_{m-i,i}^{1}\omega_{m-i,i}+M_{m-i,i-1}s_{m-i,i-1}-N_{m-i-1,i}%
s_{m-i-1,i}%
\end{align*}
for all $i$. Whence $\lim_{k}s_{i,j}^{\left(  k\right)  }=s_{i,j}$ exists for
all $i,j$, $i+j<n$, $\lim_{k}\partial_{i,j}^{1}\omega_{i,j}^{\left(  k\right)
}=\partial_{i,j}^{1}\omega_{i,j}$ and $\chi_{i,j}=\partial_{i,j}^{1}%
\omega_{i,j}+M_{i,j-1}s_{i,j-1}-N_{i-1,j}s_{i-1,j}$ hold for all $i,j$,
$i+j\leq n$.
\end{proof}

\end{document}
